% 6. Discussion (transferred from docs/paper/paper_draft.md)
\section{Discussion}
\label{sec:06_discussion}

\textbf{Graph gain without a biological gain.} In all three experiments, graph-constrained models achieved higher ROC AUC than the CNN baseline, yet the Drosophila graph did not outperform its matched controls. Mean differences went in both directions (−0.015, +0.017, −0.002) and were always small relative to their variability. The gain over the CNN therefore cannot be attributed to biological wiring. The added block (BatchNorm, two sparse nonlinear layers, a full readout) changes the architecture and adds about 2,500 parameters, and random, lattice-like and hub-dominated graphs deliver the same benefit. The graph block may be useful as an architectural component independently of which biological edges it contains.

\textbf{Why Drosophila and rewired are nearly identical.} Rewiring preserves every node's in- and out-degree: which neurons receive many inputs and which broadcast widely. It destroys partner identity, reciprocity and clustering. In our setting, image features are assigned to nodes arbitrarily (feature \emph{i} → node \emph{i}) and the encoder is learned jointly. The encoder can therefore adapt its feature-to-node assignment to whatever wiring is present. Under such co-adaptation, degree statistics plausibly matter more than exact partner identities, and the near-zero Phase 10 difference is consistent with this.

\textbf{Why the visual graph is a stronger test than EB.} The EB graph was selected by a generic synapse-count rule from a navigation centre and was dense before thresholding (0.79), which required a density-motivated threshold. The visual graph is selected by cell-type and retinotopic annotations of a well-characterised circuit, is naturally sparse (0.061) and keeps all biological edges. If any topology-specific effect of visual wiring on an image task existed at this scale, this graph offered a better opportunity to observe it. It did not.

\textbf{Negative results are informative.} Connectome-inspired architectures are sometimes motivated by the assumption that biological wiring carries a transferable computational advantage. Our results show that, at least in this small-data medical setting, apparent benefits of a connectome-wired model can be fully reproduced by matched non-biological graphs. Reporting such controlled negative results guards against attributing architectural effects to biology.

\textbf{Structural realism versus task usefulness.} The visual graph is biologically faithful: its strongest connections reproduce known motion-pathway links. Biological realism of a structure does not imply usefulness for an unrelated task (human tumour classification). Moreover, our binary mask discards synapse counts, neurotransmitter signs and cell-type identity, all of which could matter for a biologically meaningful computation.

\textbf{Small-data behaviour and variability.} With 171 patients, fold-to-fold variation (for example CNN fold-level ROC AUC ranging from 0.803 to 0.905) exceeds most between-model differences among graph models. Results on the single 21-patient test split (Phase 9) were more optimistic than cross-validated estimates (Drosophila 0.944 vs 0.925 fold-level), which illustrates the value of full-cohort evaluation.

\textbf{Calibration and thresholds.} ROC AUC is threshold-independent, whereas balanced accuracy, sensitivity and specificity depend on the frozen 0.5 threshold. At this threshold the graph models' predictions were shifted towards glioma (higher sensitivity, lower specificity). We did not test the cause, which could involve calibration, class weighting or the architecture. The lower pooled OOF AUC relative to fold-level AUC also points to calibration differences between fold models. We did not tune thresholds, to avoid test-set optimisation.

\textbf{Possible site-related confounding.} The fact that every \texttt{MR}-prefixed patient is a glioma may reflect acquisition site, so any model could partly exploit site-specific image characteristics. This affects all models equally and therefore does not bias the topology comparison, but it limits how absolute performance can be interpreted.

\textbf{No clinical implication.} The study concerns inductive biases in a controlled setting. The reported performance is not evidence of diagnostic utility, and the models are not intended for clinical use.
